\section{Conclusion}
\label{sec:conclusion}

We fine-tuned an encoder, UltraBERT-380M, as a decision model with one mask per option, and compared four training recipes.
On the same training examples, every recipe is more accurate than Qwen3-0.6B trained with LLM-as-Jev (\TestUbList{}\% to \TestUbSig{}\% against \TestQwFull{}\% on the test set).
The encoder is faster (\MsUb{} against \MsQwFull{}\,ms per decision) but uses more peak memory (\MemPeakUb{} against \MemPeakQwFull{}\,GiB).

Among the recipes, scoring by ``yes'' minus ``no'' gives the larger test gain and all of the external gain.
Sigmoid normalisation adds \DiffUbSigUbYesNoAbs{} test points, mostly on MuSiQue, and brings label frequencies closer to the gold ones; the square-root softplus does not, and we cannot say which property of the sigmoid is responsible.
An auxiliary head on the frozen encoder predicts annotated evidence in the same pass, and on the softmax encoder it ranks the annotated sentences higher than AT2 and ContextCite, except on ContractNLI.

The results have limits.
They cover one encoder, one decoder family and one English training mixture, below one billion parameters; a 4B decoder is still more accurate except on the intent datasets.
We did not evaluate whether the evidence head reflects what the model used, and we measured only peak memory.

Two directions follow: larger encoders, to see whether the advantage holds above one billion parameters, and training the decision and the evidence together, so that the evidence head can shape the encoder.
